%!TEX root = ../hutbthesis_main.tex

% 设置英文摘要
\keywordsen{Fixed-wing Aircraf; Landing Control; Human-Machine Interaction Enhancement}
\begin{abstracten}

The landing phase of a fixed-wing aircraft constitutes a critical stage
within the entire flight process---one characterized by elevated risk
and stringent control requirements---as it is susceptible to various
influencing factors such as velocity, sink rate, attitude deviations,
lateral drift errors, and terrain conditions. This paper focuses on a
fixed-wing aircraft landing control system enhanced through
human-machine interaction. The research encompasses a comprehensive
scope, including system architecture and requirements analysis, dynamic
modeling, simulation platform implementation, multi-scenario result
analysis, and the design of human-machine interaction enhancement
strategies. First, the landing mission scenarios, key parameters, and
influencing factors were thoroughly analyzed. Second, dynamic and
simulation models for the landing phase were established, followed by
the verification of a baseline operational scenario. Third, through
comparative simulations across diverse terrain conditions---including
flat, undulating, ridge, valley, and step-like profiles---the behavioral
patterns of the aircraft\textquotesingle s ground clearance, velocity,
and attitude during the final approach were analyzed; this analysis
identified the step-like terrain as a representative adverse scenario.
Finally, a human-machine interaction enhancement strategy was proposed
that integrates risk identification, status prompting, control
assistance, and constraint protection. The findings indicate that this
approach can, to a significant extent, enhance the safety and stability
of fixed-wing aircraft landings under complex operational conditions,
thereby providing a valuable reference for the subsequent optimization
of control systems.

\end{abstracten}
